\documentclass[11pt]{article}
\usepackage{amsmath,amsthm,amssymb,amsfonts}
\usepackage[margin=1in]{geometry}
\usepackage[colorlinks=true,linkcolor=blue,citecolor=blue,urlcolor=blue]{hyperref}

\title{Three Classical Analytic Estimates for Explicit-Formula Coefficients of the Riemann Zeta Function}
\author{Robert Duran\thanks{Email: \texttt{duran.robert301@gmail.com}.}}
\date{}

\theoremstyle{plain}
\newtheorem{theorem}{Theorem}[section]
\newtheorem{lemma}[theorem]{Lemma}
\newtheorem{corollary}[theorem]{Corollary}

\theoremstyle{definition}
\newtheorem{definition}[theorem]{Definition}
\newtheorem{remark}[theorem]{Remark}

\begin{document}
\maketitle

\begin{abstract}
We formulate three classical analytic estimates for coefficients arising from the explicit formula for $\psi(x)$: (i) a dyadic $L^2$ tail bound, (ii) a large-sieve--type root-mean-square inequality for spectral sums over zero ordinates, and (iii) a uniform root-mean-square tail bound for a weighted explicit-formula expansion. The inputs are entirely classical---Stirling asymptotics, bounds for $\zeta'/\zeta$ in the critical strip, and Riemann--von Mangoldt local zero-counting.
\end{abstract}

\section{Introduction}
Let $\rho = \beta + i\gamma$ range over the non-trivial zeros of the Riemann zeta function, paired under the functional equation. Write the explicit-formula contribution to $\psi(x)$ in the form
\[
  \log R(s) = \sum_{\rho} a_\rho\, e^{(\rho-1)s} + \cdots, \qquad
  j_{\mathrm{rel}}(s) = \sum_{\rho} \frac{\rho-1}{\rho}\, a_\rho\, e^{(\rho-1)s} + \cdots,
\]
and define coefficients
\[
  b_\rho := \frac{\rho-1}{\rho}\, G(\rho), \qquad
  G(s) := \pi^{-s/2} \Gamma\!\left(\frac{s}{2}\right) H(s).
\]
Here $H(s)$ is a holomorphic factor coming from Mellin transforms of test functions, rational weights, and possible $\zeta'/\zeta$ terms, assumed to be of at most polynomial growth in vertical strips with $\Re(s+\kappa)\ge 1/2+\varepsilon$ for some $\kappa,\varepsilon>0$. We present three self-contained analytic estimates for these coefficients and their dyadic aggregates. No reference is made to ERURH or to RH in what follows.

\section{Preliminaries}
We assume throughout:
\begin{itemize}
  \item Stirling-type bounds for $\Gamma(s/2)$ and $\pi^{-s/2}$ in vertical strips $\Re s \ge 1/2+\varepsilon$.
  \item Bounds for $\zeta'/\zeta(s)$ in $\Re s \ge 1/2+\varepsilon$, of the standard polynomial type in $|\Im s|$.
  \item Riemann--von Mangoldt zero-counting in dyadic intervals: for $T$ large,
  \[
    N(T) - N(T/2) \ll T \log T,
  \]
  and corresponding local counting in short ranges of length $\asymp T$.
  \item Standard estimates for exponential sums with frequencies $\gamma$ in a dyadic band (large-sieve/spacing), as used in Montgomery--Vaughan.
\end{itemize}
All implied constants may depend on $\varepsilon>0$ and on the test function parameters implicit in $H(s)$.

\section{Theorem A: Dyadic $L^2$ tail bound}
\begin{theorem}[HbWeak\_{$L^2$} tail]\label{thm:A}
Assume the classical explicit formula structure where $G(s)$ includes the gamma factor $\Gamma(s/2)$. There exist constants $C>0$, $A>0$, and $T_0>0$ such that, for all $T \ge T_0$,
\[
  \sum_{\substack{\rho \\ T < |\gamma| \le 2T}} |b_\rho|^2 \;\le\; C\, T\, (\log T)^A.
\]
\end{theorem}

\begin{proof}[Proof sketch]
Recalling that $b_\rho = \frac{\rho-1}{\rho} G(\rho)$, the presence of the factor $\Gamma(\rho/2)$ in $G(\rho)$ implies, by Stirling's formula, that $|b_\rho| \ll \exp(-\frac{\pi}{4}|\gamma|) |\gamma|^{k}$ for some fixed $k$. This exponential decay dominates any polynomial growth. Since the number of zeros in the interval is $N(2T)-N(T) \ll T \log T$, the sum is exponentially small (of order $e^{-cT}$). Therefore, for sufficiently large $T$, it is trivially bounded by $C T (\log T)^A$.
\end{proof}
\begin{remark}
The proof in fact yields an exponentially small dyadic tail in $T$; we record only a polynomial bound because it suffices for applications.
\end{remark}

\section{Theorem B: A large-sieve inequality for zero ordinates}
Fix $L>0$ and a dyadic band $T<|\gamma|\le 2T$ with $T$ large. Define the spectral sum
\[
  F(u) = \sum_{T<|\gamma|\le 2T} b_\rho\, e^{i\gamma u}.
\]

\begin{theorem}[LSGammaWeak]\label{thm:B}
Let $1 \ll L \ll T^\alpha$ with $0<\alpha<1$ (admissible range). There exist constants $C>0$, $A>0$, and $T_0>0$ such that, for $T \ge T_0$,
\[
  \frac{1}{L} \int_0^L |F(u)|^2 \, du
    \;\le\; C\,(\log T)^A \sum_{T<|\gamma|\le 2T} |b_\rho|^2.
\]
\end{theorem}

\begin{proof}[Proof sketch]
Expanding the mean square involves the kernel $K_L(t) = \frac{1}{L}\int_0^L e^{itu}\,du = e^{itL/2}\frac{\sin(tL/2)}{tL/2}$. A direct reference for exponentials with well-spaced real frequencies is Montgomery~\cite[Theorem~1.3]{MontgomeryLNM227}, which gives
\[
  \int_0^L \Bigl|\sum_\rho a_\rho e^{i\gamma_\rho u}\Bigr|^2 du \;\le\; (L + \Delta^{-1}) \sum_\rho |a_\rho|^2
\]
whenever $|\gamma_\rho - \gamma_{\rho'}| \ge \Delta$ for $\rho \neq \rho'$. In our situation, standard zero-spacing estimates in a dyadic band $[T,2T]$ allow us to choose $\Delta$ with $\Delta^{-1} \ll \log T$, so that $(L + \Delta^{-1}) \ll L + \log T$. Hence
\[
  \frac{1}{L} \int_0^L |F(u)|^2 \, du \;\ll\; \Bigl(1 + \frac{\log T}{L}\Bigr) \sum_{T<|\gamma|\le 2T} |b_\rho|^2 \;\ll\; (\log T)^A \sum_{T<|\gamma|\le 2T} |b_\rho|^2
\]
for some $A>0$ in the admissible range $1 \ll L \ll T^\alpha$.
\end{proof}

\begin{corollary}[Max-bound version]
Under the same assumptions,
\[
  \frac{1}{L} \int_0^L |F(u)|^2 \, du
    \;\le\; C'\, T \,(\log T)^{A+1} \cdot \max_{T<|\gamma|\le 2T} |b_\rho|^2.
\]
\end{corollary}
\begin{proof}
This follows immediately from Theorem \ref{thm:B} and the trivial bound $\sum |b_\rho|^2 \le (N(2T)-N(T)) \max |b_\rho|^2 \ll T (\log T) \max |b_\rho|^2$.
\end{proof}

\section{Theorem C: A uniform RMS tail bound}
Let $S$ be large, $1 \ll L \ll S^\alpha$ with $0<\alpha<1$, and fix $B<4$. Define the tail observable
\[
  \widetilde F_{\mathrm{tail}}(s) \;:=\; s^{-2} \sum_{\substack{|\gamma|>S^B}} b_\rho\, e^{i\gamma s}.
\]
  Decompose the tail region $\{|\gamma|>S^B\}$ into dyadic bands $(2^k S^B, 2^{k+1} S^B]$, $k=0,1,2,\dots$. Define
\[
  K_{\mathrm{tail}}(S,L)^2 \;:=\; \frac{1}{L} \int_{S}^{S+L} \bigl|\widetilde F_{\mathrm{tail}}(s)\bigr|^2 \, ds.
\]
\begin{theorem}[A2-tail]\label{thm:C}
Assume Theorem~\ref{thm:A} (HbWeak\_{$L^2$} tail) and Theorem~\ref{thm:B} (LSGammaWeak) on all bands in the tail. Then there exist constants $M_{\mathrm{tail}}>0$, $S_0>0$ such that, for all $S \ge S_0$ and all admissible $L$,
\[
  K_{\mathrm{tail}}(S,L)^2 \;\le\; M_{\mathrm{tail}}^2,
\]
where $M_{\mathrm{tail}}$ depends only on the constants in Theorems~\ref{thm:A} and \ref{thm:B} and on $B,\alpha$, not on $S,L$.
\end{theorem}
\begin{proof}[Proof sketch]
By the gamma factor inside $G(\rho)$ and Stirling's formula (as used in the proof of Theorem~\ref{thm:A}), $|b_\rho| \ll \exp(-c|\gamma|)|\gamma|^{k}$, so both $\sum |b_\rho|$ and $\sum |b_\rho|^2$ converge absolutely and decay like $\exp(-c S^B)$ beyond $|\gamma|>S^B$. Apply Theorem~\ref{thm:B} on each band to bound the windowed RMS contribution; summing the resulting geometric series over $k$ yields $K_{\mathrm{tail}}(S,L)^2 \ll \exp(-c' S^B)$, which is in particular $\le M_{\mathrm{tail}}^2$ for large $S$.
\end{proof}

\section{Remarks}
\begin{itemize}
  \item The proofs require only standard tools: Stirling, zero-density/spacing inputs for $\zeta'/\zeta$, Riemann--von Mangoldt, and basic large-sieve estimates for exponentials.
  \item Theorems~\ref{thm:A} and \ref{thm:B} imply exponential dyadic decay; we state polynomial forms because they suffice for many applications.
  \item Variants with test functions or different normalizations of $b_\rho$ can be handled similarly by adjusting $H(s)$ and the admissible range.
\end{itemize}

\begin{thebibliography}{9}
\bibitem{Titchmarsh} E.\,C.~Titchmarsh, \emph{The Theory of the Riemann Zeta-Function}, 2nd ed., revised by D.\,R.~Heath-Brown, Oxford University Press, 1986.
\bibitem{MontgomeryVaughan} H.\,L.~Montgomery and R.\,C.~Vaughan, \emph{Multiplicative Number Theory I.~Classical Theory}, Cambridge University Press, 2007.
\bibitem{Ivic} A.~Ivi\'c, \emph{The Riemann Zeta-Function: Theory and Applications}, Dover, 2003.
\bibitem{MontgomeryLNM227} H.\,L.~Montgomery, \emph{Topics in Multiplicative Number Theory}, Lecture Notes in Mathematics 227, Springer, 1971.
\bibitem{ZeroDensity} See, e.g., standard surveys on zero-density and zero-spacing estimates for $\zeta(s)$ for bounds on $\zeta'/\zeta$ in $\Re s \ge 1/2+\varepsilon$.
\end{thebibliography}

\end{document}
